% ==========================================
% SECTION 1: INTRODUCTION
% ==========================================
\section{Giới Thiệu Đề Tài}

% --- SLIDE 1: BỐI CẢNH NGHIÊN CỨU ---
\subsection{Bối cảnh nghiên cứu}
\begin{frame}{Bối Cảnh Nghiên Cứu \& Lý Do Chọn Đề Tài}
    \begin{block}{Giới hạn vật lý của CPU lõi đơn (Single-core)}
        \small
        Việc tăng hiệu năng bộ vi xử lý bằng cách tăng tần số xung nhịp đã đạt tới giới hạn vật lý do hiện tượng quá nhiệt và tiêu hao năng lượng cực lớn (\textbf{Power Wall, Thermal Wall}).
    \end{block}
    
    \vspace{0.15cm}
    
    \begin{block}{Chuyển dịch tất yếu sang Đa lõi (Multi-core)}
        \small
        Tích hợp nhiều lõi xử lý hoạt động độc lập trên cùng một chip để tối ưu hiệu năng tính toán song song.
    \end{block}

    \vspace{0.15cm}

    \begin{alertblock}{Bài toán đặt ra}
        \small
        Giải quyết các vấn đề phức tạp về giao tiếp liên lõi, tranh chấp bus bộ nhớ dùng chung (memory contention) và đồng bộ giữa các luồng xử lý phần cứng.
    \end{alertblock}
\end{frame}

% --- SLIDE 2: CHỦ ĐỀ DỰ ÁN ---
\subsection{Chủ đề của dự án}
\begin{frame}{Chủ Đề Của Dự Án}
    Dự án tập trung nghiên cứu, thiết kế phần cứng ở mức RTL sử dụng ngôn ngữ mô tả phần cứng \textbf{Verilog HDL} để mô phỏng bộ xử lý 8 lõi độc lập:
    
    \vspace{0.2cm}
    \begin{columns}[T]
        \begin{column}{0.48\textwidth}
            \begin{block}{Kiến trúc nhân đơn}
                \small
                Mỗi lõi CPU dựa trên kiến trúc tập lệnh mở \textbf{RISC-V} (\textbf{RV32I subset}) áp dụng kỹ thuật \textbf{pipeline 5 tầng} nhằm tối ưu hóa thông lượng lệnh thực thi.
            \end{block}
        \end{column}
        
        \begin{column}{0.48\textwidth}
            \begin{block}{Kiến trúc đa nhân}
                \small
                Cả 8 nhân CPU kết nối và sử dụng chung một khối bộ nhớ dữ liệu (\textbf{Shared Data Memory}) thông qua một bộ điều khiển Bus và phân xử tập trung.
            \end{block}
        \end{column}
    \end{columns}
\end{frame}

% --- SLIDE 3: MỤC TIÊU THIẾT KẾ ---
\subsection{Mục tiêu thiết kế}
\begin{frame}{Mục Tiêu Thiết Kế}
    Dự án đề ra ba mục tiêu cốt lõi cần đạt được:
    
    \vspace{0.25cm}
    \begin{itemize}
        \item \textbf{Thiết kế lõi đơn:} Thiết kế lõi xử lý RISC-V có pipeline 5 tầng hoàn chỉnh, xử lý xung đột lệnh (data/control hazards) một cách tối ưu bằng Forwarding, Stall và Flush.
        \item \textbf{Thiết kế trục bus đa lõi:} Xây dựng cơ chế phân xử công bằng (Round-Robin Arbiter), đảm bảo cả 8 nhân CPU đều có thể ghi/đọc bộ nhớ dùng chung mà không xảy ra xung đột hay deadlock.
        \item \textbf{Mô phỏng chức năng:} Sử dụng bộ công cụ nguồn mở \textbf{Icarus Verilog} để chạy mô phỏng hệ thống và \textbf{GTKWave} để phân tích giản đồ thời gian (waveform).
    \end{itemize}
\end{frame}